\documentclass[aspectratio=169]{beamer}

\usetheme{Madrid}
\usecolortheme{default}
\setbeamertemplate{navigation symbols}{}

\title{Reproducing and Extending a Paper on Particle Diameter Prediction}
\author{Your Name}
\institute{Your Institute}
\date{\today}

\newcommand{\imgplaceholder}[2]{
  \begin{center}
    \fbox{
      \begin{minipage}[c][#2][c]{#1}
        \centering
        \textbf{Image Placeholder}\\[0.5em]
        Replace with your figure
      \end{minipage}
    }
  \end{center}
}

\begin{document}

\begin{frame}
  \titlepage
\end{frame}

\begin{frame}{Problem Statement}
  \begin{itemize}
    \item Particle diameter is critical in biomedical micro/nano systems.
    \item Full experimental exploration is costly and time-consuming.
    \item Goal: predict particle diameter from process parameters using ML.
  \end{itemize}
  \vspace{0.3cm}
  \imgplaceholder{0.75\linewidth}{0.28\textheight}
\end{frame}

\begin{frame}{Paper Overview}
  \begin{itemize}
    \item Target: diameter prediction for chitosan and TiO2-related particles.
    \item Baseline approach: ANN models (MLP and RBF).
    \item Inputs: process parameters; Output: particle diameter.
  \end{itemize}
  \vspace{0.3cm}
  \imgplaceholder{0.75\linewidth}{0.28\textheight}
\end{frame}

\begin{frame}{What I Reproduced}
  \begin{itemize}
    \item Rebuilt data workflow and model pipeline from the paper.
    \item Implemented baseline training and evaluation steps.
    \item Compared reproduced metrics with paper-reported performance.
  \end{itemize}
  \vspace{0.2cm}
  \begin{block}{Key Message}
    The original method was implemented as closely as possible.
  \end{block}
\end{frame}

\begin{frame}{Reproduction Results}
  \begin{itemize}
    \item Baseline results follow the trend reported in the paper.
    \item Any small gap can come from split randomness and implementation details.
    \item Reproduction validates the paper's methodology in this environment.
  \end{itemize}
  \vspace{0.3cm}
  \imgplaceholder{0.85\linewidth}{0.30\textheight}
\end{frame}

\begin{frame}{My Extensions}
  \begin{itemize}
    \item Added a new model beyond the original paper.
    \item Added a new dataset to test broader applicability.
    \item Added an optimization phase for better tuning.
  \end{itemize}
  \vspace{0.3cm}
  \imgplaceholder{0.85\linewidth}{0.30\textheight}
\end{frame}

\begin{frame}{Comparison: Baseline vs Extended Pipeline}
  \begin{itemize}
    \item \textbf{Paper baseline} $\rightarrow$ reproduced.
    \item \textbf{+ New model} $\rightarrow$ improved candidate behavior.
    \item \textbf{+ New dataset} $\rightarrow$ better robustness check.
    \item \textbf{+ Optimization} $\rightarrow$ stronger final performance.
  \end{itemize}
  \vspace{0.3cm}
  \imgplaceholder{0.85\linewidth}{0.30\textheight}
\end{frame}

\begin{frame}{Conclusion}
  \begin{itemize}
    \item I reproduced the paper's core workflow and results pattern.
    \item I extended the work with a new model, new dataset, and optimization.
    \item The pipeline is both reproducible and improvable.
  \end{itemize}
  \vspace{0.3cm}
  \begin{block}{Takeaway}
    This project validates the original study and shows practical next-step improvements.
  \end{block}
\end{frame}

\begin{frame}{Thank You}
  \centering
  Questions?
\end{frame}

\end{document}
